\documentclass[11pt]{letter}

%\usepackage{pxfonts}
%\usepackage{color}
%\hyphenation{ma-nu-script}
\usepackage[colorlinks=true,urlcolor=blue, hyperindex,breaklinks=true] {hyperref}

\usepackage{graphicx}
\usepackage[a4paper,left=2cm,right=2cm,bottom=2cm]{geometry}

\usepackage{fancyhdr}
%\usepackage{marvosym}

% nice font.
\renewcommand*\rmdefault{ppl}

\usepackage[table,xcdraw,dvipsnames]{xcolor}
\newcommand{\nick}[1]{\textcolor{red}{[#1]}}
\newcommand{\ddd}[1]{\textcolor{blue}{[#1]}}

\pagestyle{empty}

\begin{document}


\hfill\begin{minipage}{6cm}
	\raggedright
	\vspace{-2.5cm}
	{\footnotesize\itshape Nikolaos Koukoulekidis}\\
	{\footnotesize\itshape Department of Physics}\\
	{\footnotesize\itshape Imperial College London}\\
	{\footnotesize\itshape London SW7 2AZ, UK}\\
	{\footnotesize\itshape %\MVAt~ 
		\verb|nk2314@imperial.ac.uk|}\\
	~\\
		\footnotesize\today
\end{minipage}

\vspace{-3.5cm}
\includegraphics[scale=.1]{icllogo.png}
\vspace{3.5cm}

\vspace{-1.5cm}
Dear Editors of PRX Quantum,

\vspace{.3cm}
	
Please find enclosed a copy of our manuscript \emph{``Constraints on magic state protocols from the statistical mechanics of Wigner negativity''}, which we would like to submit for review in PRX Quantum. 
We believe our work fits naturally into the scope of the journal. 

A fundamental problem is to understand how quantum physics admits higher degrees of order than possible in classical physics. One signature of this non-classical order is the appearance of `negative probabilities' in the Wigner representation of quantum systems. However statistical mechanics cannot be directly applied to quantify this order since the Boltzmann entropy is not well-defined when Wigner negativity arises. In this work we do two main things: we show how a well-defined statistical mechanics of Wigner negativity can be constructed, and then apply this framework to quantify the order in `magic states', which are the key sources of non-classicality in leading quantum computing models. 

It is a major goal to properly understand the processing and `distillation' of magic states for quantum technologies. Upper bounds on magic state distillation protocols do exist and are based on abstract information-theoretic measures. Our methods outperform existing upper bounds in the literature, and to our knowledge, provide the first general framework to incorporate dependence on hardware physics, such as system Hamiltonians or temperatures. As such this approach is relevant to the challenge of designing and building quantum computers.

Our results are also of interest beyond quantum technologies. We show that recent research on the thermodynamics of small quantum systems can be extended to Wigner negativity/quasi-distributions. To our knowledge, we provide the first application of majorization (a fundamental concept in classical statistics) to quasi-probability distributions, and so open up the possibility of new tools to analyse non-classicality more generally. We show that sub-set of classical (R\'{e}nyi) entropies remain meaningful on quasi-distributions and the appearance of negative entropies is a witness of non-classicality in quantum systems. Finally, our analysis also raises novel questions relevant to work on classical statistical mechanics.

For the above reasons, we believe that our manuscript is of relevance to a wide audience of quantum physicists, and so is suitable for publication in PRX Quantum.

\vspace{1cm}
\hspace{8cm}
\begin{minipage}{9cm}
\flushleft
Sincerely Yours,\\

Nikolaos Koukoulekidis, David Jennings.
\end{minipage}
\newpage


\textbf{Popular Summary}

\ddd{You can use the first paragraph, which is 100 words, for the 100 word summary. We can include the cover-letter too, which is not 100 words. Find a place to add it in and make it clear e.g. in the title, it's a cover letter. This way the editor should open it and read it.}

\nick{this is supposed to be ``three paragraphs, first should be self-contained summary within summary, second and third more extensive description'', should be at most 250 words long. Currently: 250 words}

A fundamental problem is to understand how quantum physics admits higher degrees of order than classical physics. Here we show how quantum-mechanical order can be described in terms of a rigorous statistical mechanics framework in which negative entropies arise. We then apply this framework to quantify the order in `magic states' - crucial resources that power quantum computers. Our methods outperform existing magic bounds, and provide the first general framework to incorporate dependence on hardware physics, such as system Hamiltonians or temperatures. This work is relevant both to quantum technologies and the fundamental challenge of quantifying non-classicality in quantum systems.

A signature of non-classical order is the appearance of `negative probabilities' in representations of quantum systems. However traditional statistical mechanics cannot be directly applied to negative probabilities. We first show how to extend statistical mechanics by using a more fundamental concept of order (`majorization') from statistical theory and then apply this to quantifying the order in magic states.

Quantum computers `burn' magic states as a form of fuel to provide exponential speed-ups for important tasks, such as in computational Chemistry. Our results show that this thermodynamic intuition can be made precise and that magic states, in a rigorous sense, carry `negative entropy' that is processed thermodynamically. Our methods outperform existing magic bounds in the literature, and to our knowledge, provide the first general framework to incorporate dependence on hardware physics, such as system Hamiltonians or temperatures. As such this approach is relevant to the challenge of designing and building quantum computers.




\end{document}
